% !TEX root = main.tex
\subsection{Related Work}\label{subsection:introduction-related-work}

\medskip
\noindent\textbf{Signability in prior notions.}
Table~\ref{table:comparisons-with-other-privacy-preserving-signatures} summarizes the landscape.
Here, we discuss the work most closely related to \SMS in design goals.
We classify \MPS~\cite{CRYPTO:NGSY22} as limited function composition:
the signer evaluates a public function to obtain an index $j$,
then selects the $j$-th member of a public family $\{P_i\}_{i = 1}^k$
as the effective hidden policy.
The hidden function is never directly composed with a public one;
the family size is bounded ($k = 4$ in~\cite{CRYPTO:NGSY22});
and the selection circuit evades, rather than solves,
the challenge of proving hidden-circuit satisfiability.
\SMS replaces selection with genuine composition via the combiner $P$,
supports an unbounded family $\FFamily$ of hidden arithmetic circuits,
and directly confronts hidden-circuit evaluation.
Furthermore, \MPS's selection-based approach costs
$\BigO(k \cdot (\nmult + \nadd))$ constraints;
the PLONK-based reduction developed here reduces this to $\BigO(\nmult + \nadd)$,
removing the multiplicative factor $k$.
The original PBS construction~\cite{PKC:BelFuc14},
built on Groth--Sahai proofs~\cite{EC:GroSah08}
over structure-preserving signatures~\cite{CRYPTO:AFGHO10},
is inherently restricted to functions expressible as bilinear maps.
Subsequent lattice-based constructions~\cite{ChengNW16,XuSNW22}
extend the reach only to linear systems over $\ZZ_2$.
Bobolz~\etal~\cite{FC:BobDiaKoh24} introduced universal anonymous signatures
to balance utility and privacy,
but omit hidden signing functions
and provide only a high-level framework without concrete instantiations.
Nguyen~\etal~\cite{AC:NRSX23} defined bicameral and auditably private signatures,
pairing a private signing policy with a public disclosure mechanism
in a fixed two-state model;
that binary structure is a special case of \SMS,
which supports arbitrary composition of hidden and public functions
and a continuously variable accountability spectrum.

\medskip
\noindent\textbf{Accountability mechanisms in prior notions.}
The accountability axis has a richer history than the binary GS model suggests.
Linkable ring signatures, proposed by Liu, Wei, and Wong~\cite{ACISP:LiuWeiWong04}, allow linking two signatures
by the same signer without revealing identity.
Accountable ring signatures, introduced by Xu and Yung~\cite{CARDIS:XuYun04}
and first realized in the standard model by
Lai, Zhang, Chow, and Schr\"{o}der~\cite{ESORICS:LaZhChowSch16},
allow the signer to designate an opener who can later reveal the signer's identity,
thereby combining the ad hoc group flexibility of ring signatures
with the accountability of group signatures.
Chow, Susilo, and Yuen~\cite{VIETCRYPT:ChowSusYue06} proposed
escrowed linkability, where only a designated linking authority
can decide whether two ring signatures share a signer.
Traceable ring signatures, proposed by Fujisaki and Suzuki~\cite{PKC:FujSuz07}, 
publicly reveal any signer who signs two distinct messages under the same tag, without any designated opener.

Kiayias, Tsiounis, and Yung~\cite{EC:KiayiasYung04} introduced traceable signatures,
where a tracing trapdoor tags a misbehaving user's signatures without opening others;
Chow~\cite{SAC:Chow09} improved this by enabling per-user tagging
without first collecting all signatures.
Abe, Chow, Haralambiev, and Ohkubo~\cite{IJIS:AbeChowHarOhk13}
introduced double-trapdoor anonymous tags for a modular construction of traceable signatures
supporting authorship claiming and denial:
the signer herself can selectively prove or disprove her authorship of a specific signature,
a form of user-controlled accountability orthogonal to authority-initiated tracing.
Kiayias and Zhou~\cite{FC:KiayiasZhou07} proposed hidden identity-based signatures,
where the opening authority needs only its secret key to reveal the signer's identity;
Chow, Zhang, and Zhang~\cite{FC:ChowZhaZha17} gave the first standard-model construction.
Zhang, Wu, and Chow~\cite{RSA:ZhaWuChow19} introduced certified limited opening,
delegating opening rights so that neither the master certifier
nor unauthorized parties can revoke anonymity.
Each of these offers a distinct tracing granularity;
\SMS's opening tag $\op$ can encode any of these levels
as a special configuration of the combiner $P$;
user-controlled accountability such as authorship claiming and denial
points to a natural extension of the framework.

\medskip
\noindent\textbf{Proofs for hidden function evaluations.}
Several recent works construct zero-knowledge proofs
for the satisfiability of hidden circuits,
but each either targets a different security model or a different problem class.
Ling~\etal~\cite{PQCRYPTO:LNPTW21} gave zero-knowledge proofs
for committed symmetric Boolean functions,
covering only a narrow subclass.
Boneh, Nguyen, and Ozdemir~\cite{EPRINT:BonNguOzd21} commit to
a hidden circuit's PLONK's description via Marlin~\cite{EC:CHMMVW20},
yielding proofs of arithmetic-circuit satisfiability,
but their construction relies on the algebraic group model (AGM)~\cite{CRYPTO:FucKilLos18},
which carries documented soundness concerns~\cite{AC:ZhaZhoKat22}.
Di~\etal~\cite{AC:DXNT24} propose MUXProofs,
layering Marlin over a univariate sum-check to verify RAM steps encoded as rank-1 constraint systems;
Choudhuri~\etal~\cite{PETS:CGGSS24} introduced Sublonk,
applying PLONK's arithmetization to arbitrary-degree polynomials
for instruction-level proofs.
Both also rely on the AGM.

Our construction operates in the random oracle model (ROM) for both
the generic and the lattice-based instantiation, avoiding these concerns entirely.
Tang, Pham, and Ngo~\cite{EPRINT:TanPhaNgo24} apply PLONK's arithmetization
to prove correct execution of RAM programs,
a problem structurally different from that of hidden-circuit evaluation for signatures;
the treatment developed here, specifically treating the PLONK's description
as a witness component to reduce hidden-circuit satisfiability
to a fixed public verification circuit,
is tailored to the NIZKAoK requirements of privacy-preserving signatures
and is not addressed in that prior work.
In short, no prior scheme achieves ROM-based proofs
for arbitrary hidden bounded-size arithmetic circuits
in a privacy-preserving signature setting.

\medskip
\noindent\textbf{Relationship to anonymous credentials.}
Anonymous credentials~\cite{SJ:AuSMC13,EC:CamLys01,CRYPTO:CamLys04} and \SMS share structural ancestry:
both involve a certifying authority, a hidden user attribute,
and a zero-knowledge proof of possession.
The two paradigms are nonetheless orthogonal.
The former focuses on selective disclosure of attribute \emph{values}
in a presentation token, where the predicate being proven is public to the verifier.
\SMS focuses on signing messages under a hidden function \emph{structure}~$F$,
whose circuit topology is concealed from signatures,
and adds per-signature accountability control absent from anonymous credentials,
which typically target full unlinkability.
Combining the two, producing an \SMS signature while selectively disclosing
which attributes the signer holds, is a natural direction for future work
and is part of the broader functional credentials framework~\cite{iotpts/Chow16}.